\section{Introduction to $\infty$-categories}
\textbf{Speaker: Malthe Sporring}

\subsection{Introduction}
\subsubsection{What are $\infty$-categories?}
\begin{itemize}
    \item Objects and $n$-morphisms for each $n$, which satisfy associativity and commutativity (weakyl up to a higher morphism).
    \item \textbf{The homotopy hypothesis:} $\infty$-groupoids are weak homotopy types. This was historically known to be true for $n$-groupoids for low $n$. Here, a  homotopy type is a topological space up to weak homotopy equivalence. A weak homotopy equivalence is a map $f:X\rightarrow Y$ which induces bijections $\pi_n(X)\rightarrow \pi_n(Y)$ for each $n$.
    \item We take the homotopy hypothesis as a definition of $\infty$-groupoids. We then identify $\infty$-categories as "categories enriched in homotopy types". In particular, for each pair of objects $(x,y$), we ask for a homotopy type $Map(x,y)$ of maps from $x$ to $y$.
\end{itemize}

\subsubsection{Why are $\infty$-categories needed?}
Consider the $1-$category $hTop:=Top[w.h.e.^{-1}]$. Its objects are topological spaces up to weak homotopy equivalences, and its morphisms are zig-zags $X\rightarrow Y\leftarrow Y'\rightarrow Z\dots \rightarrow B$ where only weak homotopy equivalences are allowed to go backwards.

We would like a notion of homotopy colimit. At the very least this should be invariant under replacing a diagram by a weakly homotopy equivalent one. However, the following suggests that it will not be enough to take colimits in $hTop$:

\begin{remark}
    The span $\bullet \leftarrow S^0\rightarrow \bullet$ in $Top$ has pushout $\bullet$. On the other hand, the span $D^1 \leftarrow S^0\rightarrow D^1$ has pushout $S^1.$ However there is a weak homotopy equivalence  $D^1s \simeq \bullet$.
\end{remark}

Indeed, it turns out $hTop$ doesn't have colimits at all.

Instead, we would like homotopy colimits to have a universal property which is weaker than that of a colimit. In particular, let $A\xleftarrow{f} B\xrightarrow{g} C$ be a span in a category $\mathcal{C}$. Then a homotopy pushout should be an object $D$, morphisms $j:A\rightarrow D$ and $h:C\rightarrow D$, \textit{along with a homotopy} $hg\Rightarrow jf$. Its universal property should state that given 

There is no way to state the full universal property within an $n$-category. For the full universal property we need $\mathcal{C}$ to be an $\infty$-category.

There is a practical, and somewhat ad-hoc, solution to this (given in the setting of a model category), which we discuss briefly below. However, it doesn't satisfy a familiar universal property. In general, the motto is the following: \textbf{homotopy-coherent analogues of universal properties in category theory only exist in the language of $\infty$-categories.}

\subsubsection{Suggested references}
\begin{itemize}
    \item (Start here) Synthetic approach: \cite{Haugseng-categories-intro}
    \item Quasi-categories: \cite{Land}
    \item For details / reference: \cite{HTT}
\end{itemize}

\subsection{Models of $\infty$-categories}
\subsubsection{Model categories}
A model category is a category $C$ along with wide subcategories $(W,Fib,Cofib)$ of \textit{weak equivalences, fibrations,} and \textit{cofibrations}. In a model category, one can identify those diagrams whose ordinary colimit corresponds to homotopy colimits, and prove that any diagram is weakly equivalent to such a well-behaved diagram. For instance, there is a model structure on $Top$ where the weak equivalences are the weak homotopy equivalences.

Model categories can be very useful for computing specific homotopy colimits. However, they are inconvenient for a systematic generalization of ordinary category theory for a number of reasons:
\begin{itemize}
    \item As previously mentioned, familiar universal properties of homotopy colimits, etc, do not exist in a model category.
    \item There is no (large) model category of (small) model categories. This is a massive problem for 
    \item More generally, categories like $Ab$ have universal properties among (presentable, additive, symmetric monoidal) categories. While one can identify model categories (e.g. of spectra) which are analogous to $Ab$, one cannot identify the corresponding universal property in a model category.
    \item In general, there are many model categories which present the same $\infty$-category. When working with model categories one ends up repeatedly having to prove that two model categories are equivalent (Quillen equivalent), which is (a) hard, and (b) Quillen equivalences don't behave as well as one would expect. For instance, a Quillen equivalence $C\rightarrow D$ does not imply the existence of a Quillen equivalence $D\rightarrow C$. Furthermore, if $D\leftarrow C\rightarrow E$ is a span of Quillen equivalences, there is no requirement of the existence of a Quillen equivalence $D\rightarrow E$. On the other hand, any zig-zag of Quillen equivalences automatically gives rise to an equivalence of underlying $\infty$-categories!
\end{itemize}

For these reasons, $\infty$-categories are much better suited as a setting to do formal higher category theory and higher algebra! 

Furthermore, there is a wealth of results which assures homotopy theorists that the "standard" notions of, e.g., homotopy colimit, localise to the correct notions of colimit in $\infty$-categories. Therefore, one can always descend to a model category, for instance to do some computation, and then safely transport the result back to the $\infty$-categorical level.

\subsubsection{Quasi-categories}
The most developed model of $\infty$-categories, and the one used in \cite{HA}, is that of quasi-categories. These were invented by Boardmann-Vogt, developed by Joyal and later Lurie. The landmark textbook \cite{HTT} of Lurie remains the most extensive reference for $\infty-$category theory. Quasi-categories are called $\infty$-categories in $\cite{HTT}$ and in $\cite{HA}$ (our main reference for this reading group).

Let $\Delta$ be the category of non-empty finite ordered sets and order-preserving maps. We define  $sSet:=Fun(\Delta^{op},Set)$.

\begin{remark}
    Think of $X_0$ as the set of objects, $X_1$ as the set of morphisms, $X_2$ as the set of compositions, etc. Think of the simplicial maps as giving $id:X_0\rightarrow X_1$, $s,t:X_1\rightarrow X_0$ and $\circ:X_2\rightarrow X_1$. 
\end{remark}

\begin{remark}
Think of a homotopy $f\Rightarrow g$ as witnessing a composition $g=id\circ f$.     
\end{remark}

We let $\Delta^n$ be the simplicial set $[n]$. We let the $i^{th}$ horn $\Lambda^{n}_i$ be obtained from $\Delta^n$ by removing the $i'th$ face. 

Note that we have functors $\Delta^{op}\rightarrow Cat$ and $\rightarrow Top$. Left Kan extension along the Yoneda embedding gives functors $|-|:sSet\rightarrow Top$ and $h:sSet \rightarrow Cat$, left adjoint to the functors $Sing:Top\rightarrow sSet$ and $N:Cat\rightarrow sSet$ defined by $N(C)_n=Fun([n],C)$ and $Sing(X)_n=Fun(\Delta^n_{Top},X)$. In this way we can view categories, and topological spaces, as simplicial sets.

\begin{definition}

\begin{itemize}
    \item A Kan fibration is a map of simplicial sets $f:X\rightarrow Y$ with the right lifting property with respect to all horn fillers.
    \item A Kan complex is a simplicial set whose unique map to $\Delta^0$ is a Kan fibration.
    \item An inner fibration is a map with the RLP with respect to inner horn fillers.
    \item A quasicategory is a simplicial set whose unique map to $\Delta^0$ is an inner fibration.
\end{itemize}    
\end{definition}

We have a notion of homotopy (and hence homotopy groups etc) internal to simplicial sets. A simplicial homotopy is a map $S\times \Delta^1\rightarrow T$. A map $f$ of simplicial sets is a weak homotopy equivalence if and only if $f$ is a weak homotopy equivalence. In particular, we say a simplicial set $X$ is (weakly) contractible if the map $X\rightarrow \Delta^0$ is a weak homotopy equivalence.

Note that for a quasicategory, we are not requiring unique composites. However, the simplicial set of composites is a contractible Kan complex.

We let the category $Kan\subset sSet$ be the full subcategory spanned by the Kan complexes. The category of simplicial sets is naturally enriched over itself by letting, for each Kan complex $X_\bullet,Y_\bullet$, the simplicial set $$Fun(X_\bullet,Y_\bullet)_{n}:=sSet(X_\bullet\times\Delta^n,Y_\bullet).$$ The category $Kan$ is also naturally self-enriched.

There is a natural way to consider $[n]$ as a simplicial category. Left Kan extension of $\Delta\rightarrow Cat_{sSet}$ along the Yoneda embedding $y:\Delta\rightarrow sSet$ once again provides a left adjoint $sSet\rightarrow Cat_{sSet}$, whose right adjoint is called the \textbf{coherent nerve}. The coherent nerve of the simplicial category $Kan$ is the quasi-category of Kan complexes, which is a model for $\infty$-groupoids.

One can identify the quasi-category of quasi-categories in a similar way. A map of quasicategories is a categorical equivalence if it is "fully faithful and essentially surjective". These need to be forced to be equivalences. However, as the next section shows, it may be enough to stop at a model of $\infty$-groupoids...

\subsubsection{Complete Segal Spaces}
Categories are categories internal to $Set$. Suppose we have a model $\mathcal{S}$ for $\infty$-groupoids (e.g. Kan complexes), and a model for homotopy pullbacks, then define $\infty$-categories as "$\infty$-categories internal to $\mathcal{S}$".

\begin{definition}
    A Segal space is a functor $\Delta^{op}\rightarrow \mathcal{S}$ satisfying the Segal condition:
    \begin{itemize}
        \item The natural map $X_{n}\rightarrow  X_1\times_{X_0}\dots \times_{X_0} X_1$ is an equivalence. (Here we take the homotopy pull-back)
    \end{itemize}
    it is furthermore \textbf{complete}, if equivalences in $X_0$ are the same as "equivalences induced by $X_1$".
\end{definition}

The latter condition hints at the following: $X_0$ has a notion of what it means for two objects $x$ and $y$ to be equivalent: just the existence of an equivalence $x\rightarrow y$. On the other hand, one can ask for the existence of arrows $x\rightarrow y$ and $y\rightarrow x$, as objects in $X_1$, along with 2-simplices, as an object in $X_2$, witnessing $g\circ f\simeq id$ and $f\circ g \simeq id$. It is not automatic that these two notions of equivalence agree, so we need to specify it.

Complete Segal spaces model $\infty$-categories. An advantage of this model is that you can complete all your calculations in $\mathcal{S}.$ For instance, colimits in functor $\infty$-categories are computed pointwise.

Another advantage is that it is often easier to check that a category is a complete Segal space, as opposed to a quasi-category.

\subsubsection{The synthetic approach}
In practice, one can eventually forget the model once it has proven a few key things. One can state these key assertions, and leave it to the reader to either take them as axiomatic or prove in their favourite model that they hold. Afterwards we can all forget what model we chose.

\subsection{Assertions}
These are (some of) the assertions of \cite{Haugseng-categories-intro}, and are not meant to be formal. In general, the rest of this text is greatly based on \cite{Haugseng-categories-intro}.

\begin{itemize}
    \item There exist $\infty$-groupoids, also called spaces. These have points, paths of points, homotopies of paths, etc. A map is an equivalence if there exist homotopies, etc. For each pair of points $x,y$ in an $\infty$-groupoid, there exists the mapping space $Map(x,y)$.
    \item There is an identity morphism $id_x:Map(x,x)$. There exist composition functors $Map(x,y)\times Map(y,z)\rightarrow Map(x,z)$. Composition and associativity is unital up to associativity.
    \item Sets are $\infty$-groupoids. $\emptyset$ is the initial $\infty$-groupoid, and $\{*\}$ is the terminal $\infty$-groupoid.
\end{itemize}
\begin{definition}
    We say a space is (weakly) contractible if the unique map to $*$ is an equivalence.
\end{definition}
\begin{definition}
    We say a commutative square
\[\begin{tikzcd}
	x & z \\
	y & w
	\arrow["g"', from=1-1, to=1-2]
	\arrow["f", from=1-1, to=2-1]
	\arrow["h"', from=1-2, to=2-2]
	\arrow["j", from=2-1, to=2-2]
\end{tikzcd}\]
is the data of the $1-$morphisms depicted and a $2$-morphism $j\circ f\Rightarrow h\circ g$.
    
\end{definition}
\begin{itemize}
    \item for an $\infty$-groupoid, there exist its set $\pi_0(X)$ of path components.
    \item (Homotopy) pull-back squares of $\infty$-groupoids exist. The points of $X\times_Z Y$ are points of $X$ and $Y$ and a path between them in $Z.$
\end{itemize}

\begin{definition}
    For every space $X$ and point $x\in X$,  we define $\Omega_x(X)$ to be the pull-back

\[\begin{tikzcd}
	\Omega(X) & x \\
	x & {X}
	\arrow[from=1-1, to=1-2]
	\arrow[from=1-1, to=2-1]
	\arrow[from=1-2, to=2-2]
	\arrow[from=2-1, to=2-2]
	\arrow["\ulcorner"{anchor=center, pos=0.125, rotate=180}, draw=none, from=2-2, to=1-1]
\end{tikzcd}\]

    We furthermore define $\pi_n(X,x)=\pi_0(\Omega_x^n(X))$.
\end{definition}

\begin{itemize}
    \item $\pi_n(X,x)$ is a group for $n\geq 1$, and abelian for $n\geq 2.$
    \item Homotopy groups detect equivalences, i.e. a map is an equivalence if and only if it induces an isomorphism on $\pi_i$ for all $i$.

\end{itemize}

\begin{proposition}\label{prop:equiv}
    A map $f:X\rightarrow Y$ is an equivalence if and only if it induces a surjection $f:\pi_0(X)\rightarrow \pi_0(Y)$ and equivalences $X(x,y)\rightarrow Y(f(x),f(y))$ for each $x,y\in X$.
\end{proposition}
\begin{proof}
    Exercise for the reader. Hint: show that the conditions automatically imply that $f:\pi_0(X)\rightarrow \pi_0(Y)$ is injective.
\end{proof}

\begin{itemize}
    \item There exist $\infty-$categories, which form an $\infty$-category, functors of them, and the category of functors is an $\infty$-category. (Homotopy) pull-backs and pushouts of $\infty$-categories exist.
    \item Ordinary categories are $\infty$-categories. 
\end{itemize}



\begin{itemize}
    \item For each $\infty$-category $C$, the $\infty$-groupoids $C^\simeq$ and $|C|$ exist. $C^\simeq $ is the space $Map([0],C)$.
    \item A functor $F:C\rightarrow D$ is an equivalence if and only if $C^\simeq \rightarrow D^\simeq$ and $Map([1],C)\rightarrow Map([1],D)$ are equivalences of spaces.
    \item Opposite $\infty$-categories exist.
    \item The degeneracy map $id_-:C^\simeq\rightarrow Ar(C)$ is fully faithful. 
\end{itemize}

An equivalence in an $\infty$-category $C$ is an arrow in $C^\simeq$. Equivalently, a map $f:x\rightarrow y$ is an equivalence if, for every $a\in C$, we have that the induced functors
$$Map(a,x)\rightarrow Map(a,y)$$ and $$Map(y,a)\rightarrow Map(x,a)$$ are equivalences.


\subsection{Basic constructions}
We can make many definitions using the following definition:
\begin{definition}
    Let $f:x\rightarrow y$ and $g:z\rightarrow w$ be morphisms. We say $f$ is left orthogonal to $g$ (equivalently, $g$ is right orthogonal to $f$) if, for each commutative square as below, the $\infty$-groupoid of fillers

\[\begin{tikzcd}
	x & z \\
	y & w
	\arrow[from=1-1, to=1-2]
	\arrow[from=1-1, to=2-1]
	\arrow[from=1-2, to=2-2]
	\arrow[dashed, from=2-1, to=1-2]
	\arrow[from=2-1, to=2-2]
\end{tikzcd}\]
is contractible.
    
\end{definition}

We write $[n]$ for the ordinary category $\sf 1\rightarrow 2\rightarrow \dots \rightarrow n$.

\begin{definition}
    We say a functor $F:C\rightarrow D$ of $\infty$-categories is conservative if it is right orthogonal to the unique functor $[1]\rightarrow [0]$.
\end{definition}

In ordinary category theory, a functor $F:C\rightarrow D$ is conservative exactly if an arrow $f$ in $C$ is an isomorphism only if $F(f)$ is an isomorphism. Convince yourself that this agrees with the above definition if $F$ is a functor of ordinary categories. \textbf{Warning:} even though all the functors involved are ordinary functors, don't forget that commutativity only holds up to a $2$-morphism! In the case of ordinary categories, $2$-morphisms correspond to natural isomorphisms.

We let $\delta[1]$ be the maximal subgroupoid of $[1]$. I.e. it is the category with two objects $\sf \{1,2\}$ and no non-trivial morphisms.


\begin{definition}
    We say a functor $F:C\rightarrow D$ of $\infty$-categories is fully faithful if it is right orthogonal to the canonical functor $\delta[1] \rightarrow [1]$.
\end{definition}

In ordinary category theory, a functor $F:C\rightarrow D$ is fully faithful exactly if, for each $x,y\in C$ it induces a bijection $C(x,y)\rightarrow D(F(x),F(y)$. Convince yourself that this agrees with the above definition in ordinary category theory.

\begin{proposition}\label{prop:FF}
    A functor $F:C\rightarrow D$ is fully faithful if and only if, for each $x,y\in C$, the induced functors $Map(x,y)\rightarrow Map(F(x),F(y))$ are equivalences.
\end{proposition}

\begin{theorem}\label{thm:FF-is-cons}
    Fully faithful functors are conservative.
\end{theorem}
\begin{proof}
    See \cite{Haugseng-categories-intro}[Proof of 2.7.1]
\end{proof}

\begin{definition}
    A functor $F:C\rightarrow D$ is essentially surjective if it induces a surjection $F:\pi_0(C^\simeq)\rightarrow \pi_0(D^\simeq)$.
\end{definition}

%We now get the following, which is sometimes baked into the definition of an $\infty$-functor, as a result.


\begin{proposition}
    Conservative functors are preserved under pullback.
\end{proposition}

For each $x,y\in C$, we can identify the mapping space $Map(x,y)$ space as a pull-back

% https://q.uiver.app/#q=WzAsNCxbMCwwLCJNYXAoeCx5KSJdLFsxLDAsIkFyKEMpIl0sWzEsMSwiQ1xcdGltZXMgQyJdLFswLDEsIlxce3gseVxcfSJdLFszLDJdLFsxLDJdLFswLDNdLFswLDFdLFswLDIsIiIsMSx7InN0eWxlIjp7Im5hbWUiOiJjb3JuZXIifX1dXQ==
\[\begin{tikzcd}
	{Map(x,y)} & {Ar(C)} \\
	{\{x,y\}} & {C\times C}
	\arrow[from=1-1, to=1-2]
	\arrow[from=1-1, to=2-1]
	\arrow["\lrcorner"{anchor=center, pos=0.125}, draw=none, from=1-1, to=2-2]
	\arrow[from=1-2, to=2-2]
	\arrow[from=2-1, to=2-2]
\end{tikzcd}\]

The functor $(s,t):Ar(C)\rightarrow C\times C$ is conservative, so it follows that $Map(x,y)\rightarrow \{x,y\}$ is conservative. This implies $Map(x,y)$ is an $\infty$-groupoid.

\begin{proposition}
    A conservative functor $F:C\rightarrow D$ induces a monomorphism $F:\pi_0(C^\simeq)\rightarrow \pi_0(D^\simeq)$.
\end{proposition}

\begin{prop}
    A functor is an equivalence if and only if it is fully faithful and a surjection on $\pi_0$.
\end{prop}

\subsection{coCartesian fibrations and straightening-unstraightening}
The following is an $\infty$-categorical version of the Grothendieck construction of ordinary category theory. It is extremely important for $\infty$-category theory.

\subsubsection{Left/right fibrations}

\begin{definition}
    A functor $\pi:E\rightarrow B$ is a left fibration if it is right orthogonal to the functor $[0]\rightarrow [1]$ given by $0\mapsto 0$. It is a right fibration if it is right orthogonal to the functor $[0]\rightarrow [1]$ where $0\mapsto 1.$

    We let the category of left fibrations over $B$, $left_{/B}\subset Cat_{\infty/B}$, be the full subcategory spanned by the left fibrations. We similarly define $right_{/B}$.
\end{definition}

\begin{example}\label{ex:left-fib}
    Given a category $C$ and an object $x\in C$, define the slice category $C_{/x}$ by the pull-back 
\[\begin{tikzcd}
	{C_{/x}} & {Ar(C)} & C \\
	{\{x\}} & C
	\arrow[from=1-1, to=1-2]
	\arrow[from=1-1, to=2-1]
	\arrow["\lrcorner"{anchor=center, pos=0.125}, draw=none, from=1-1, to=2-2]
	\arrow["s"', from=1-2, to=1-3]
	\arrow["t"', from=1-2, to=2-2]
	\arrow[from=2-1, to=2-2]
\end{tikzcd}\]
Then the functor $s:C_{/x}\rightarrow C$ is a right fibration. Similarly, the functor $t:C_{x/}\rightarrow C$ is a left fibration. Instructive exercise !
\end{example}

\begin{prop}
    Left/right fibrations are preserved by pull-back.
\end{prop}

\begin{example}
    Consider Example \ref{ex:left-fib} in the case $C=Gpd_{\infty}$. and $x=\Delta^0$. The category $(Gpd_{\infty})_{\Delta^0/}$ is the category of pointed spaces, and left fibration $t:(Grpd_\infty)_{\Delta^0/}\rightarrow (Grpd_\infty)$ forgets the pointing.

    Given a functor $F:B\rightarrow Gpd_\infty$, we can consider the left fibration $\int F\rightarrow B$ given by pull-back along $t:(Grpd_\infty)_{\Delta^0/}\rightarrow (Grpd_\infty)$. We call $\int F\rightarrow B$ the \textbf{unstraightening} of $F$.
\end{example}


\begin{theorem}[Straightening/unstraightening I]
 (Covariant/contravariant) unstraightening gives equivalences $Fun(B,Gpd_{\infty})\rightarrow left_{/B}$ and $Fun(B^{op},Gpd_\infty)\rightarrow right_{/B}$.
\end{theorem}
We call the inverse the straightening functor. It takes a left fibration $E\rightarrow B$ to a functor $B\rightarrow Gpd_{\infty}$ with the property that $b\mapsto E_b$ (up to equivalence). Given a morphism $b\rightarrow b'$, the induced functor $F:E_b\rightarrow E_{b'}$ sends an object $x\in E_b$ to the codomain of the unique lift 
\[\begin{tikzcd}
	{[0]} & E \\
	{[1]} & B
	\arrow["x", from=1-1, to=1-2]
	\arrow[from=1-1, to=2-1]
	\arrow[from=1-2, to=2-2]
	\arrow[dotted, from=2-1, to=1-2]
	\arrow["{b\rightarrow b'}"', from=2-1, to=2-2]
\end{tikzcd}\]

Note that this choice $x\mapsto F(x)$ is only defined up to contractible choice. This is a feature, not a bug! To give a functor $F:B\rightarrow Gpd_\infty$ is to make choices of spaces $F(b)$, functors $F(b)\rightarrow F(b')$, $2-$cells $\dots.$ On the other hand, to prove that a functor $E\rightarrow B$ is a left fibration is to check a \textbf{property}, not to give a \textbf{structure}. Think of a left fibration $E\rightarrow B$ as having \textbf{the ability} to choose a functor $B\rightarrow Gpd_\infty$, uniquely up to contractible choice, \textbf{without ever making any of those choices explicit.}

\begin{exercise}
    Prove that given a left fibration $E\rightarrow B$ all the fibers $E_b, b\in B$ are $\infty$-groupoids.
\end{exercise}

\subsubsection{coCartesian fibrations}

\begin{definition}
    Let $\pi:E\rightarrow B$ be a functor. An edge $f:x\rightarrow y$ of $E$ is \textbf{$\pi$-coCartesian} if, for each other $z\in E$ the diagram given by precomposition with $f$,

    % https://q.uiver.app/#q=WzAsNCxbMCwwLCJFKHkseikiXSxbMSwwLCJFKHgseikiXSxbMCwxLCJCKFxccGkoeSksXFxwaSh6KSkiXSxbMSwxLCJCKFxccGkoeCksXFxwaSh6KSkiXSxbMCwxLCJmXyEiXSxbMiwzLCJcXHBpKGYpXyEiLDJdLFswLDIsIlxccGkiLDJdLFsxLDMsIlxccGkiXV0=
\[\begin{tikzcd}
	{E(y,z)} & {E(x,z)} \\
	{B(\pi(y),\pi(z))} & {B(\pi(x),\pi(z))}
	\arrow["{f_!}", from=1-1, to=1-2]
	\arrow["\pi"', from=1-1, to=2-1]
	\arrow["\pi", from=1-2, to=2-2]
	\arrow["{\pi(f)_!}"', from=2-1, to=2-2]
\end{tikzcd}\]

    is a pull-back diagram.
\end{definition}

\begin{definition}
    Let $\pi:E\rightarrow B$ be a functor. An edge $f:x\rightarrow y$ of $E$ is \textbf{$\pi$-Cartesian} if, for each other $z\in E$ the diagram given by post-composition with $f$,

% https://q.uiver.app/#q=WzAsNCxbMCwwLCJFKHoseCkiXSxbMSwwLCJFKHoseSkiXSxbMCwxLCJCKFxccGkoeiksXFxwaSh4KSkiXSxbMSwxLCJCKFxccGkoeiksXFxwaSh5KSkiXSxbMCwxLCJmXyoiXSxbMiwzLCJcXHBpKGYpXyoiLDJdLFswLDIsIlxccGkiLDJdLFsxLDMsIlxccGkiXV0=
\[\begin{tikzcd}
	{E(z,x)} & {E(z,y)} \\
	{B(\pi(z),\pi(x))} & {B(\pi(z),\pi(y))}
	\arrow["{f_*}", from=1-1, to=1-2]
	\arrow["\pi"', from=1-1, to=2-1]
	\arrow["\pi", from=1-2, to=2-2]
	\arrow["{\pi(f)_*}"', from=2-1, to=2-2]
\end{tikzcd}\]

    is a pull-back diagram.
\end{definition}

\begin{definition}
A functor $\pi: E\rightarrow B$ is a \textbf{coCartesian fibration} if , for each $x\in E$ and $f:\pi(x)\rightarrow b'$ in $B$, there exists a $\pi-$coCartesian morphism $x\rightarrow y$ lifting $f$.

A functor $\pi: E\rightarrow B$ is a \textbf{Cartesian fibration} if , for each $y\in E$ and $f:b\rightarrow \pi(y)$ in $B$, there exists a $\pi-$coCartesian morphism $x\rightarrow y$ lifting $f$.
\end{definition}

\begin{example}
    Given a category $C$, the target functor $t:Ar(C)\rightarrow C$ is a coCartesian fibration. The source functor $s:Ar(C)\rightarrow C$ is a Cartesian fibration. An edge is $s$-Cartesian if and only if it is inverted by $t$, and is $t$-coCartesian if and only if it is inverted by $s.$ Instructive exercise!
\end{example}

\begin{exercise}
    Prove that $\pi$-coCartesian lifts of morphisms are unique up to equivalence.
\end{exercise}

\begin{exercise}
    Prove that for a coCartesian fibration $\pi:E\rightarrow B$, the composite of two $\pi$-coCartesian edges is $\pi-$coCartesian.
\end{exercise}

\begin{exercise}
Prove that left fibrations are coCartesian fibrations. Prove that a coCartesian fibration $\pi:E\rightarrow B$ is a left fibration if and only if every edge is $\pi$-coCartesian.
\end{exercise}

\begin{prop}
    coCartesian fibrations are preserved by pull-back.
\end{prop}

\begin{definition}
We let $coCart_{/B}\subset Cat_{\infty/B}$ be the (non-full) subcategory on the coCartesian fibrations and functors that preserve coCartesian edges.
\end{definition}


\begin{theorem}[Straightening-Unstraightening]
There is a functor $coCart_{/B}\rightarrow Fun(B,Cat)$, called straightening, which is an equivalence. Its inverse is called (covariant) unstraightening.

Similarly, there is functor $Cart_{/B}\rightarrow Fun(B^{op},Cat)$, called straightening, which is an equivalence. Its inverse is called (contravariant) unstraightening.
\end{theorem}

Straightening takes a coCartesian fibration $\pi:E\rightarrow B$ to the functor $F:B\rightarrow Cat$ which sends $b\in B$ to the category $E_b$ and $b\rightarrow b'$ to the functor $E_b\rightarrow E_b'$ which sends an object $x\in E_b$ to the codomain of the unique $\pi-$coCartesian morphism $x\rightarrow y$ over $b$.

Since $\pi$-coCartesian lifts are only unique up to contractible choice, the choice $x\mapsto F(x)$ is also only defined up to contractible choice. This is a feature, not a bug! To give a functor $F:B\rightarrow Cat_\infty$ is to make choices of categories $F(b)$, functors $F(b)\rightarrow F(b')$, $2-$cells $\dots.$ On the other hand, to prove that a functor $E\rightarrow B$ is a coCartesian fibration is to check a \textbf{property}, not to give a \textbf{structure}. Think of a coCartesian fibration $E\rightarrow B$ as having \textbf{the ability} to choose a functor $B\rightarrow Cat_\infty$, uniquely up to contractible choice, \textbf{without ever making any of those choices explicit.}

\subsubsection{Adjunctions}

An adjunction is a functor $C\rightarrow \Delta^1$ which is both a coCartesian fibration and a Cartesian fibration.

The covariant straightening produces a functor $\Delta^1\rightarrow Cat_\infty$, which is the same as a functor $F:C_0\rightarrow C_1$ in $Cat_\infty$. The contravariant unstraightening produces a functor $G:C_1\rightarrow C_0$. Furthermore, $F$ is left adjoint to $G$! As the next exercise shows, this can be reduced to familiar definitions.

\begin{exercise}
    Convince yourself that a coCartesian and Cartesian fibration $C\rightarrow \Delta^1$ as below encodes unit and counit natural transformations $\eta:id_{C_0}\Rightarrow GF$ and $\epsilon: FG\Rightarrow id_{C_1}$ by identifying natural components of such natural transformations for each $x\in C_0$ and each $y\in C_1$.
\end{exercise}


\subsection{Factorisation systems}
The replacement for factorisation systems in ordinary category theory are orthogonal factorisation systems. These are also just called factorisation systems.

\begin{definition}
    An orthogonal factorisation system on a category $C$ is two full subcategories $C_L,C_R$, such that 
    \begin{itemize}
        \item $C_L\cap C_R\simeq C^\simeq$,
        \item Every map $f:x\rightarrow y$ in $C$ has a factorisation as a map in $C_L$ followed by a map in $C_R$,
        \item Maps in $C_L$ are left orthogonal to maps in $C_R$.
    \end{itemize}
\end{definition}

This definition is equivalent to asking for the space $Fact(f)$ of factorisations of $f$ as a map in $C_L$ followed by a map in $C_R$ is contractible. This is non-obvious, and it is highly instructive exercise to work out why these definitions are equivalent.

\begin{proposition}
    Let $\pi:E\rightarrow B$ be a Cartesian fibration. Let $\pi-Cart\subset E$ be the wide subcategory of Cartesian morphisms, and $\pi-inv\subset E$ the subcategory of morphisms sent to equivalences by $\pi.$ Then $(\pi-inv,\pi-Cart)$ is a factorisation system on $E$.

    Similarly, if $\pi':E'\rightarrow B'$ is a coCartesian fibration, then $(\pi-coCart,\pi-inv)$ is a factorisation system on $E'$.
\end{proposition}
\begin{proof}
    Instructive exercise!
\end{proof}

\begin{example}
    There is a factorisation system on $Ar(C)$ given by $s-$inv morphisms followed by $t$-inverted morphisms. This arises from the Cartesian fibration $s:Ar(C)\rightarrow C$, after noting $t-inv=s-Cart$.
\end{example}

\subsection{Representables}
\begin{definition}
    A representable fibration is a right fibration $s:C_{/x}\rightarrow C$.
\end{definition}

Note that the unstraightening of such a fibration is of the form $Map(-,x):C\rightarrow Gpd_\infty$.

Recall that $(s,t):Ar(C)\rightarrow C\times C$ has the property that the projection $t:Ar(C)\rightarrow C$ is coCartesian with $t$-coCart $=$ $s-inv$, and the projection $s:Ar(C)\rightarrow C$ is Cartesian with $s-Cart = t-inv.$ Such a functor is called a bifibration.

\begin{exercise}
    Convince yourself that given a bifibration $(s,t):E\rightarrow C\times C$, and an object $x\in C$, that the induced functor $s_x:E_x\rightarrow C$, given by pull-back along $C\times \{x\}\rightarrow C\times C$, is a right fibration.
\end{exercise}

The above exercise suggests the existence of a modified straightening-unstraightening equivalence, which proves an equivalence between bifibrations $(s,t):E\rightarrow C\times C$ and functors $C\rightarrow right_{/C}$ of the form $x\mapsto s_x:E_x\rightarrow C$.

Indeed this exists! Applied to the bifibration $(s,t):Ar(C)\rightarrow C\times C$ induces a functor $y:C\rightarrow RFib(C)$. This is the $\infty$-categorical version of the Yoneda embedding.

\begin{proposition}
    The Yoneda embedding is fully faithful.
\end{proposition}

\subsection{(Co)limits}
\begin{definition}
    An object $x\in C$ is terminal if, for all $a\in C$, the space $C(a,x)$ is contractible. It is initial if $C(x,a)$ is contractible.
\end{definition}

Given a functor $p:J\rightarrow C$ we can define the $\infty-$category of cocones as the pull-back in $Cat_\infty$

% https://q.uiver.app/#q=WzAsNCxbMCwwLCJDX3svcH0iXSxbMCwxLCJDIl0sWzEsMSwiRnVuKEosQykiXSxbMSwwLCJGdW4oSixDKV97cC99Il0sWzAsMV0sWzEsMiwiXFxEZWx0YV9DIl0sWzAsM10sWzMsMl0sWzAsMiwiIiwxLHsic3R5bGUiOnsibmFtZSI6ImNvcm5lciJ9fV1d
\[\begin{tikzcd}
	{C_{p/}} & {Fun(J,C)_{p/}} \\
	C & {Fun(J,C)}
	\arrow[from=1-1, to=1-2]
	\arrow[from=1-1, to=2-1]
	\arrow["\lrcorner"{anchor=center, pos=0.125}, draw=none, from=1-1, to=2-2]
	\arrow[from=1-2, to=2-2]
	\arrow["{\Delta_C}", from=2-1, to=2-2]
\end{tikzcd}\]

\begin{definition}
    A colimit of $p:J\rightarrow C$ is an initial object of $C_{p/}$.
\end{definition}

The colimit, in this sense, has the full universal property of the homotopy colimit hinted at in the introduction.